
\chapter{Study}

The major activity in this chapter is literature survey, investigation
and arriving at the targetted scope and specifications. Here the 'what',
'why' and 'how' of the product and product sub-parts are studied.
One looks for similars in already existing products and makes comparisons.
One identifies potential decisions that need to be taken with respect
to the product evolution and performs a comparative statement of possible
options that will affect design decisions. 

NOTE: While discussing the various options and the comparisions, you
need to cite appropriate references at appropriate places in the text.
A typical \emph{citation \cite{breedveld1991sic}} is referred as
shown here. The cited references will automatically be reflected in
the bibliography section in the order of occurence of citation in
the manuscript. A sample \emph{report.bib} file is included in this
template. Kindly go through it and make your bibliography file accordingly.
A free software like \emph{JABREF} (a java application) or \emph{ZOTERO}
(a firefox addon) \cite{karmakar1990dec} can be used to create your
bibliography file.

\section{Functional concept}

This should include 
\begin{itemize}
\item a discussion on the product with the aid of block diagrams and figures
\cite{thoma1976ibg}
\item an exhaustive list of stimuli to product
\item classification of the stimuli as input/output, power signal/information
signal, electric/mechanical etc.
\item First level quantification of the stimuli like, voltage levels, type
of signals (dc, ac, pulsed), rise time, current levels, etc.
\end{itemize}

\section{Module study}

Partition the entire system into modules.

\subsection{Modules}

For each of the modules, include
\begin{itemize}
\item an exhaustive list of stimuli to the module
\item classification of the stimuli as input/output, power signal/information
signal, electric/mechanical etc.
\item First level quantification of the stimuli like, voltage levels, type
of signals (dc, ac, pulsed), rise time, current levels, etc.
\end{itemize}

\subsection{Sub-modules}

If necessary, partition the various modules into sub-modules and so
on till the component level. 

\emph{NOTE}: At every module, sub-module or component level (electric
hardware or algorithms or industrial design), if there are multiple
options available, then perform a comparative study of the most significant
features among the options to evaluate the merits and demerits of
the various options. Based on the merits and demerits of the various
options, give your recommendations for design.

\section{Industrial design}

The industrial design/product design aspects should consider the following,
\begin{itemize}
\item Design mix - functional, ergonomics and aesthetics aspects of the
product should be discussed giving consideration to appropriateness
to a given application.
\item Layout - the layout of the modules and the sub-modules within an enclosure
or system should be discussed as determined by the following considerations,

\begin{itemize}
\item user, form, ergonomics, aesthetics, interconnections, display, etc.
\item realisation issues of a particular layout
\item functional aspects like EMI, thermal, wiring etc.
\end{itemize}
\item Product structure - provide a discussion on the product generic structure,
whether it should be

\begin{itemize}
\item rack and frame mounted 
\item modular cabinet
\item non-modular
\item custom 
\end{itemize}
\item Materials - provide a discussion on the choice of materials for the
product
\item Manufacturing processes
\item Cost of the product
\end{itemize}

\section{Target scope}

The following questions need to be addressed:
\begin{itemize}
\item What is the targetted scope of the product?
\item What is the set of target specifications to design the product?
\end{itemize}

